\documentclass[aspectratio=169,10pt]{beamer}
\usetheme[numbering=fraction,progressbar=frametitle]{metropolis}
\usepackage{fontspec}
\setsansfont{Latin Modern Sans}
\usepackage{amsmath,amssymb,mathtools,graphicx,booktabs}
\usepackage{polyglossia}
\setdefaultlanguage{english}
\setbeamertemplate{navigation symbols}{}
\setbeamerfont{frametitle}{size=\large}
\setbeamerfont{normal text}{size=\normalsize}
\setlength{\parskip}{.45em}
\hypersetup{pdfauthor={},pdfsubject={Lecture slides},pdfkeywords={},pdftitle={Proportions and percentage change}}
\providecommand{\N}{\mathbb{N}}
\providecommand{\Z}{\mathbb{Z}}
\providecommand{\D}{\mathbb{D}}
\providecommand{\Q}{\mathbb{Q}}
\providecommand{\R}{\mathbb{R}}
\providecommand{\C}{\mathbb{C}}
\providecommand{\vect}[1]{\overrightarrow{#1}}
\providecommand{\vi}{\vec{\imath}}
\providecommand{\vj}{\vec{\jmath}}
\providecommand{\vk}{\vec{k}}
\providecommand{\norm}[1]{\left\lVert #1\right\rVert}
\providecommand{\vc}[2]{\begin{pmatrix}#1\\#2\end{pmatrix}}
\providecommand{\deter}[2]{\det\left(#1,#2\right)}
\providecommand{\abs}[1]{\left\lvert #1\right\rvert}
\providecommand{\dd}{\,\mathrm{d}}
\providecommand{\e}{\mathrm{e}}
\providecommand{\Card}{\operatorname{Card}}
\providecommand{\E}{\operatorname{E}}
\providecommand{\V}{\operatorname{V}}
\providecommand{\degre}{^{\circ}}
\providecommand{\un}[1]{\,\mathrm{#1}}
\providecommand{\sci}[2]{#1\times10^{#2}}
\providecommand{\pgcd}{\operatorname{PGCD}}
\providecommand{\ppcm}{\operatorname{PPCM}}
\title{Proportions and percentage change}
\subtitle{Lecture and worked examples}
\date{}
\author{}
\begin{document}
\begin{frame}[plain]
\vfill
{\usebeamerfont{title}\inserttitle\par}
\medskip
{\usebeamerfont{subtitle}\insertsubtitle\par}
\medskip
{\usebeamercolor[fg]{progress bar}\rule{\textwidth}{.5pt}}
\vfill
\end{frame}
\begin{frame}{Proportion and proportion of a proportion}
In a nonempty finite set \(E\), the proportion of a subset \(F\) is \(p=\Card(F)/\Card(E)\); \(0\leq p\leq1\). Its expression as a percentage is \(100p\,\%\).
If \(G\subseteq F\subseteq E\), \(F\ne\varnothing\), and if \(p_1\) is the proportion of \(F\) in \(E\) and \(p_2\) that of \(G\) in \(F\), the proportion of \(G\) in \(E\) is \(p_1p_2\).
\end{frame}
\begin{frame}{Proportion and proportion of a proportion}
A school has \(1500\) students, of whom \(652\) are day boarders. Calculate their proportion. In a class, \(60\%\) of the students play sport; among them, \(40\%\) are boys. What proportion of the class do the boys who play sport represent?
\end{frame}
\begin{frame}{Worked answer: Proportion and proportion of a proportion}
\(652/1500\approx0.4347\), that is about \(43.47\%\). For the class, \(0.60\times0.40=0.24\), that is \(24\%\). The second percentage has the sporty students as its reference, not the whole class.
\end{frame}
\begin{frame}{Proportion and proportion of a proportion}
In a fictional school of \(800\) students, \(35\%\) play a sport at a club; among them, \(45\%\) play a team sport. Calculate the two counts and the final proportion.
\end{frame}
\begin{frame}{Worked answer: Proportion and proportion of a proportion}
\(800\times0.35=280\) students at a club. \(280\times0.45=126\) in a team sport. The final proportion is \(126/800=0.1575=15.75\%\).
\end{frame}
\begin{frame}{Proportion and proportion of a proportion}
Proportion and rate of change do not answer the same question: a proportion compares a part to a whole, a rate compares a change to an initial value.
\end{frame}
\begin{frame}{Absolute change and rate of change}
For an initial value \(V_I>0\) and a final value \(V_F>0\), the absolute change is \(V_F-V_I\); the rate of change is
\[t=\dfrac{V_F-V_I}{V_I}.\]
The rate has no unit. It is positive for an increase, negative for a decrease. The reference is always the initial value.
\end{frame}
\begin{frame}{Absolute change and rate of change}
A cheese is sold for \(2.13\) euros instead of \(3.55\) euros. Calculate the absolute change and the rate.
\end{frame}
\begin{frame}{Worked answer: Absolute change and rate of change}
Absolute change: \(2.13-3.55=-1.42\) euro. Rate: \(-1.42/3.55=-0.40\), that is a decrease of \(40\%\).
\end{frame}
\begin{frame}{Absolute change and rate of change}
\(CM=V_F/V_I=1+t\), so \(V_F=V_I(1+t)\). To recover \(V_I\), divide \(V_F\) by \(1+t\) (nonzero).
\end{frame}
\begin{frame}{Absolute change and rate of change}
A fictional item priced \(50000\) Ar increases by \(12\%\), then separately consider a \(12\%\) discount on the initial price. Give the final prices. A price after a \(20\%\) increase is \(72000\) Ar; find its initial price.
\end{frame}
\begin{frame}{Worked answer: Absolute change and rate of change}
Increase: \(50000\times1.12=56000\) Ar. Discount: \(50000\times0.88=44000\) Ar. Last question: \(V_I=72000/1.20=60000\) Ar.
\end{frame}
\begin{frame}{Successive changes}
Successive changes with rates \(t_1,t_2\) have overall multiplier
\[CM_G=(1+t_1)(1+t_2),\qquad t_G=(1+t_1)(1+t_2)-1.\]
We multiply the multipliers; we do not simply add the rates.
\end{frame}
\begin{frame}{Successive changes}
A manufacturer sold \(5600\) bicycles in 2022, which is \(25\%\) more than in 2021. Its sales decrease by \(15\%\) in 2023. Determine the sales of 2021 and 2023, then the overall rate between these two years.
\end{frame}
\begin{frame}{Worked answer: Successive changes}
2021: \(5600/1.25=4480\) bicycles. 2023: \(5600\times0.85=4760\) bicycles. Overall multiplier \(1.25\times0.85=1.0625\), so an overall increase of \(6.25\%\). Check: \((4760-4480)/4480=0.0625\).
\end{frame}
\begin{frame}{Successive changes}
A price increases by \(10\%\) then decreases by \(10\%\). Does it return to the initial price?
\end{frame}
\begin{frame}{Worked answer: Successive changes}
\(CM_G=1.10\times0.90=0.99\). The final price is \(99\%\) of the initial price: an overall decrease of \(1\%\). The two rates refer to different reference values.
\end{frame}
\begin{frame}{Successive changes}
A rate goes from \(20\%\) to \(25\%\). Explain the difference between the increase in percentage points and the relative change.
\end{frame}
\begin{frame}{Worked answer: Successive changes}
The gap is \(5\) percentage points. The relative change is \((0.25-0.20)/0.20=0.25\), that is an increase of \(25\%\).
\end{frame}
\begin{frame}{Reverse change}
To return to the initial value after a change of rate \(t>-1\), use the multiplier \(1/(1+t)\). The reverse rate is
\[t_{\text{reverse}}=\dfrac1{1+t}-1=\dfrac{-t}{1+t}.\]
\end{frame}
\begin{frame}{Reverse change}
Which rate compensates a \(25\%\) increase? Which rate compensates a \(20\%\) decrease?
\end{frame}
\begin{frame}{Worked answer: Reverse change}
After a \(25\%\) increase, the reverse multiplier is \(1/1.25=0.8\): a \(20\%\) decrease. After a \(20\%\) decrease, it is \(1/0.8=1.25\): a \(25\%\) increase.
\end{frame}
\begin{frame}{Reverse change}
A price undergoes a \(30\%\) discount, then a \(10\%\) increase. What additional percentage increase restores the initial price?
\end{frame}
\begin{frame}{Worked answer: Reverse change}
The current multiplier is \(0.70\times1.10=0.77\). The required multiplier is \(1/0.77=100/77\). The rate is \(23/77\approx0.2987\), that is about \(29.87\%\). Check \(0.77\times(100/77)=1\).
\end{frame}
\end{document}
